\begin{table}[htbp]
\centering\small
\caption{Model architectures by experiment.}
\label{tab:architectures}
\begin{tabularx}{\linewidth}{@{}>{\raggedright\arraybackslash}p{.26\linewidth}X@{}}
\toprule
Task & Model \\
\midrule
Digit classifier & Fixed multilayer perceptron with two 64-neuron tanh layers. Train a linear classifier within a fixed parameter subspace. \\
Synthetic generators and forecasting & FORCE-trained recurrent networks with 1,000 neurons for phase recovery and 256 for the periodic-signal and forecast tasks \\
CIFAR-10 interventions & Convolutional neural network (CNN) with ReLU activations and $3\times3$ kernels. Channel counts 16, 32, and 32. Max pooling after the first two layers, then adaptive average pooling and a linear classifier. \\
ResNet transfer & ResNet-18 with a ten-class classifier. Random initialization uses a $3\times3$ input convolution with stride 1 and no max pooling. ImageNet initialization uses a $7\times7$ input convolution with stride 2 and max pooling. \\
Text VAE & 48-dimensional shared embeddings, GRU encoder and decoder with hidden size 64, 16-dimensional Gaussian code \\
Grokking & One-layer transformer, width 128, four 32-dimensional heads, MLP width 512, no layer normalization; full-batch comparison uses a width-500 quadratic perceptron \\
MNIST representation selection & 784--256--256 encoder with batch normalization and ReLU; 256--256--64 projector (one variant outputs 16 coordinates); predictor width 64 in the standard SimSiam/BYOL settings \\
MNIST intervention rules & 196--128--128--10 ReLU MLP; pruning narrows both hidden layers \\
Loss-spike warning & Pre-layer-normalized transformer with two layers, width 64, and four heads \\
Walker2d & Separate 64--64 tanh MLPs for actor and critic; six action outputs \\
\bottomrule
\end{tabularx}
\end{table}

\begin{table}[htbp]
\centering\small
\caption{Classifier intervention and detector settings.}
\label{tab:cnn-protocol}
\begin{tabularx}{\linewidth}{@{}>{\raggedright\arraybackslash}p{.28\linewidth}X@{}}
\toprule
Component & Setting \\
\midrule
\multicolumn{2}{@{}l}{\textbf{CIFAR-10 intervention benchmark}} \\
Small CNN training & 10,000 CIFAR-10 images; 10,000 SGD updates; learning rate 0.02, momentum 0.9, weight decay $5\cdot10^{-4}$, batch 64 with replacement \\
Intervention update & 4,000; pre/post summaries exclude windows that overlap the intervention \\
Interventions & Divide the learning rate by 3, 10, or 100. Freeze the first two convolutions, all layers before the classifier, or all weights except the classifier bias. Prune 50\%, 80\%, or 95\% of weights. \\
Comparison records & Unchanged training or four times the batch size. Two additional records multiply the later part of an unchanged log by 10 or smooth it with a mean over the latest 16 samples. \\
Replication & 4 development and 4 held-out repeats; 10 additional repeats for timing-shift transfer \\
\addlinespace
\multicolumn{2}{@{}l}{\textbf{Detector calibration and comparisons}} \\
Detector & $M=2$, $B=3$, $D_j<-0.08829884$; feature windows ending at or before update 1,500 are excluded \\
Scoring & 5,000 updates after intervention. Each evaluation repeat includes 7 strong intervention types and 4 comparison conditions. \cx{The strong set excludes the threefold learning-rate reduction and freezing the first two convolutions.} \\
Calibration choices & $M\in\{2,3,4\}$, $B\in\{3,4,6\}$. The threshold adds a margin of 0.02 beyond the largest score in the development comparison records. \\
IAAFT surrogates & 2 reordered records per window, 100 iterations. Trimming to match the beginning and end can remove up to 15\% of the window. \\
Entropy features & Permutation order 5, delay 1, stable tie sorting, \cx{normalization by $\log(5!)$}. Sample-entropy template length 2, delay 1, Chebyshev distance, tolerance 0.2 window SDs, no self-matches. \\
CUSUM calibration & Estimate the initial mean and scale on windows ending at updates 2,000--3,500. Use a minimum scale $\max(0.01|\mu_0|,10^{-10})$. Test drift parameters $\{0.25,0.5,1\}$ with threshold margin 0.5. \\
Loss of active channels & Increase the learning rate by a factor of 20, 30, 50, or 100 for one update, 15 for three updates, or 20 for two. Measure inactivity on 500 images \cx{every 250 updates. Require the median inactive fraction in updates 5,000--9,000 to exceed that in $[2{,}000,4{,}000)$ by at least 0.05. Training accuracy below 20\% denotes network-wide failure.} \\
\cx{Training-loss comparison} & \cx{Use the median minibatch loss in each window. With $M=2$ and $B=3$, alarm if $D_j>-0.0590403$.} \\
\addlinespace
\multicolumn{2}{@{}l}{\textbf{Transfer to ResNet-18}} \\
ResNet data and SGD & Full CIFAR-10. Random initialization uses $32\times32$ inputs. ImageNet initialization uses $64\times64$ with ImageNet normalization. Batch 128 without augmentation, learning rate 0.02, momentum 0.9, weight decay $5\cdot10^{-4}$. \\
ResNet update covariance & 3 repeats per initialization type. Each update is compressed to 1,024 coordinates before computing covariance PR. \\
\bottomrule
\end{tabularx}
\end{table}

\begin{table}[htbp]
\centering\small
\caption{Text VAE penalty comparisons and check scheduling.}
\label{tab:vae-protocol}
\begin{tabularx}{\linewidth}{@{}>{\raggedright\arraybackslash}p{.28\linewidth}X@{}}
\toprule
Component & Setting \\
\midrule
Data & 6,000 training and 256 evaluation sentences, 8--24 words before the end-of-sentence marker. Vocabulary of 2,000 tokens. \\
Training & The decoder receives the actual preceding words. Adam learning rate $10^{-3}$, 32 sentences per batch, gradient norm limited to 5, zero dropout. \\
\addlinespace
\multicolumn{2}{@{}l}{\textbf{Penalty comparison}} \\
Penalty conditions & Train to update 1,024 at $\beta=0.01$. Continue from this training state to update 3,072 for each setting. Use $\beta=0.01$, $\beta=1$, or $\beta=1$ with free bits (0.5 nats per coordinate). \\
Recorded reconstruction loss & 16 fixed evaluation sentences per update. Early windows end at updates 512--1,024 and late windows at 2,048--3,072. \\
Code-use measurements & Every 128 updates. For MI, draw 4 codes per sentence from its encoder distribution. For the code-swap test, average the divergence of token predictions in both directions. \\
Comparing code use within a repeat & Loss of code use requires both measurements to fall below half of their low-penalty values. Preservation requires each to retain at least half its early value and double its value under suppression. \\
\addlinespace
\multicolumn{2}{@{}l}{\textbf{Scheduling information checks}} \\
Repeated penalty schedule & After 1,024 initial updates \cx{at $\beta=0.01$}, repeat a 2,048-update schedule three times. Raise the penalty weight $\beta$ from 0 to 1 \cx{over the first 1,024 updates}, then hold it at 1. \\
Loss and recovery labels & Measure MI and code-swap divergence every 64 updates. Divide each by its median from checks 128 updates apart between 512 and 1,024. Loss requires both ratios $\le0.50$ and recovery both $\ge0.65$ for two consecutive checks. \cx{The second check sets the event time. Score events with a full 512 updates left for detection.} \\
\cx{Eligible starting state} & \cx{The early MI and code-swap medians must exceed $0.1$ and $10^{-4}$ nats, respectively. At update 1,024 both ratios must be at least $0.65$.} \\
Check schedules & One initial check; extra budgets 6, 12, 24 \cx{with 12 primary}; at least 128 updates between triggered checks; 1 development and 9 evaluation repeats \\
\cx{Threshold grid} & \cx{$0.025,0.05,0.075,0.1,0.15,0.2,0.3,0.4,0.5,0.7,1,1.5,2,3$. Maximize detections for each method and budget. Break ties by fewer unproductive checks, shorter mean delay, fewer total checks, and a larger threshold, in that order.} \\
\bottomrule
\end{tabularx}
\end{table}

\begin{table}[htbp]
\centering\small
\caption{Walker2d training and evaluation settings.}
\label{tab:walker-protocol}
\begin{tabularx}{\linewidth}{@{}>{\raggedright\arraybackslash}p{.27\linewidth}X@{}}
\toprule
Component & Setting \\
\midrule
\multicolumn{2}{@{}l}{\textbf{Policy training}} \\
Proximal policy optimization (PPO) & Collect 512 steps in each of 8 environments per training round. Minibatch 256, 5 optimization epochs. \\
Optimizer & Learning rate decreases from $3\cdot10^{-5}$ to $3\cdot10^{-6}$. PPO clipping parameter 0.1, KL stopping parameter 0.01, gradient norm limit 0.5. \\
Training objective & Reward discount 0.99, generalized advantage estimation parameter 0.95, entropy weight 0, value-loss weight 0.5. The normalization statistics are fixed. \\
Continued training (fine-tuning) & For each penalty strength, start from the policy saved at 655,360 simulator steps and continue for 1,048,576 steps, summed across environments. \\
\addlinespace
\multicolumn{2}{@{}l}{\textbf{Command smoothing and state repeatability}} \\
Fixed-policy motion records & Discard 512 initial steps and record 4,096. The reference gait uses 152 samples. Its prescribed stride lasts about 152.215 simulator steps. \\
\cx{Reference-gait penalty} & \cx{Compare $\lambda=0$ with $\lambda=3$ at $\sigma^2=38.4774$. The reference advances by $(1+\sqrt{2}/1000)^{-1}$ samples per simulator step. Both policies resume from the same training state.} \\
State and performance criteria & In the reference-gait test, require return error and phase-matched variation ratios $\le0.75$ against training without the reference reward. Require at least eight common complete simulations, at least 90\% of the comparison reward, and at most one extra fall. Baseline return error and variation must be at least 0.02 and 0.001. \\
\addlinespace
\multicolumn{2}{@{}l}{\textbf{Policy choice under actuator shifts}} \\
Candidate policies & Use penalty strengths $\lambda\in\{0,0.25,1,4\}$ and save after 262,144, 524,288, 786,432, and 1,048,576 further steps. Four evaluation repeats. \\
Choice in unperturbed trials & Three episodes per candidate. Discard 256 steps and measure 2,048. Require at least two completed episodes and at least 90\% of the best mean reward in these trials. \\
Fixed policy choice & Use the final policy with $\lambda=0.25$, chosen on a preliminary run, if eligible. Otherwise select the highest unperturbed reward. \\
Perturbed trials & Mix in the previous filtered value with weight $\alpha\in\{0.05,0.1,0.2\}$. Use Gaussian noise of standard deviation $\sigma\in\{0,0.02\}$, with five episodes per condition. \\
\bottomrule
\end{tabularx}
\end{table}

\begin{table}[htbp]
\centering\small
\caption{Grokking runs and the matched-window analysis.}
\label{tab:grokking-protocol}
\begin{tabularx}{\linewidth}{@{}>{\raggedright\arraybackslash}p{.26\linewidth}X@{}}
\toprule
Component & Setting \\
\midrule
Modular addition & Modulus 113; training fraction 0.30; weight decay 1.0; 20,000 optimizer updates; 3 regularized runs \\
Permutation composition & $S_5$; training fraction 0.50; weight decay 0.2; 30,000 optimizer updates \cx{over 15,000 loop iterations with two updates per iteration}; 1 regularized run \\
Transformer optimization & AdamW at $10^{-3}$, \cx{$(\beta_1,\beta_2)=(0.9,0.98)$}, batch 256, float64 \\
Comparison runs & One run with zero weight decay per task. The perceptron uses half the modulo-97 addition pairs and full-batch gradient descent on squared error. \\
Matched windows & 60 log samples centered on each covariance estimate. $E\in\{4,6,10\}$, $\tau\in\{1,2,4\}$, $k\in\{5,20\}$. Exclusion uses the span or autocorrelation time. \\
\bottomrule
\end{tabularx}
\end{table}
